\topic{投票悖论}
\index{voting paradox|(}
%\emph{(Optional material from this chapter is discussed here,
%but is not required to complete the exercises.)}

想象一个研究美国总统选举进程的政治学课堂
举行一次模拟选举。
班上的 $29$ 名学生要对民主党、
共和党和第三党的候选人
从最偏好到最不偏好进行排序
（$>$ 表示「偏好于」）。
\begin{center}
  \renewcommand{\arraystretch}{1.1}
  \begin{tabular}{cc}
    \multicolumn{1}{c}{\begin{tabular}[b]{@{}c@{}}
       \textit{偏好顺序}
    \end{tabular}}
    &\multicolumn{1}{l}{{\renewcommand{\arraystretch}{1.0}
      \begin{tabular}[b]{@{}l@{}}
         \textit{具有该偏好}    \\[-0.5ex]
         \textit{的人数}
      \end{tabular}}}                \\ 
    \hline
    \begin{tabular}{@{}r@{}}
      $\text{民主党} > \text{共和党} > \text{第三党}$      \\
      $\text{民主党} > \text{第三党} > \text{共和党}$      \\
      $\text{共和党} > \text{民主党} > \text{第三党}$      \\
      $\text{共和党} > \text{第三党} > \text{民主党}$      \\
      $\text{第三党} > \text{民主党} > \text{共和党}$      \\
      $\text{第三党} > \text{共和党} > \text{民主党}$   
    \end{tabular}
    &\begin{tabular}{@{}r@{}}                                                
      $5$     \\
      $4$     \\
      $2$     \\
      $8$     \\
      $8$     \\
      $2$    
    \end{tabular}    
  \end{tabular}
    \renewcommand{\arraystretch}{1}
\end{center}
整个群体的偏好是什么样的？

总体来看，
该群体更偏好民主党而非共和党，多出五票；
把民主党排在共和党之上的选民有十七人，
反过来排的有十二人。
而该群体
更偏好共和党而非第三党候选人，
票数为十五比十四。
但奇怪的是，该群体也
更偏好第三党而非民主党，票数为十八比十一。
\begin{center}
  \includegraphics{vs/mp/voting.1}
\end{center}
这就是一个\definend{投票悖论}\index{voting paradox!definition}，
具体地说，是一个
\definend{多数循环}。\index{voting paradox!majority cycle}

数学家研究投票悖论，部分原因是它对现实政治有影响。
例如，授课教师可以操纵这个班级，
让它先在共和党与第三党之间进行一次投票，
然后在那一轮较量的胜者（将是共和党）
与民主党之间再进行一次投票，
从而选出民主党作为总体赢家。
通过类似的操纵，教师也可以让另外两位候选人中的任何一个胜出。
（我们只讨论三位候选人的选举，但在更大规模的选举中同样会出现这种情形。）

数学家研究投票悖论也因为它们本身就有趣。
一个有趣的方面是，
尽管每个选民的偏好列表都是
\definend{理性的}\index{voting paradox!rational preference order}，
即成一条直线式的顺序，
整个群体却出现了多数循环。
也就是说，多数循环似乎是在总体层面上产生的，
而并未出现在构成总体的各个成分，即偏好列表之中。
然而我们可以用线性代数论证：
循环偏好的倾向实际上存在于每个选民的列表之中，
并且当这种倾向被累加多于被抵消时
它就会显露出来。

为此，把选项缩写为 $D$、$R$ 和 $T$，
我们可以描述偏好顺序为 $D>R>T$ 的选民
对上述循环的贡献方式。
\begin{center}
  \includegraphics{vs/mp/voting.2}
\end{center}
（这里的负号是因为箭头把 $T$ 描述为偏好于 $D$，
而这位选民的偏好恰好相反。）
其他偏好列表的描述见\pageref{table:Voting} 页的表格。

现在，为了进行这次选举，我们对这些描述作线性组合；
例如，政治学课堂的模拟选举
\begin{equation*}
  5\cdot \votinggraphic{43}
  +4\cdot\votinggraphic{44}
  +\dots+2\cdot\votinggraphic{45}
\tag*{}\end{equation*}
就产生了前面展示过的环形群体偏好。

当然，取线性组合正是线性代数。
这种图形化的循环记法很有启发性但不便使用，
于是我们改用列向量：从 $D$ 出发，
按逆时针顺序依次取循环中的数字。
这样，我们把模拟选举和一张 $D>R>T$ 的选票表示如下。
\begin{equation*}
  \votepreflist{7}{1}{5}
  \quad\text{和}\quad
  \votepreflist{-1}{1}{1}
\tag*{}\end{equation*}

我们将把选票向量分解成两部分，
一部分是循环的，另一部分是非循环的。
对于第一部分，
我们称一个向量为\definend{纯循环的}，如果它属于
$\Re^3$ 的这个子空间。
\begin{equation*}
  C
  =\set{\votepreflist{k}{k}{k} \suchthat k \in\Re}
  =\set{k\cdot\votepreflist{1}{1}{1} \suchthat k \in\Re}
\end{equation*}
对于第二部分，考虑与 $C$ 中所有向量
都垂直的向量所成的集合。
\nearbyexercise{exer:PerpIsSubSp} 表明这是一个子空间
\begin{align*}
  C^\perp &=\set{\votepreflist{c_1}{c_2}{c_3} \suchthat 
                 \colvec{c_1 \\ c_2 \\ c_3}\dotprod\colvec{k \\ k \\ k} =0
                 \text{\ 对所有 $k\in\Re$}}             \\
          &=\set{\votepreflist{c_1}{c_2}{c_3} \suchthat 
                 c_1+c_2+c_3=0}                                         
          =\set{c_2\votepreflist{-1}{1}{0}
                +c_3\votepreflist{-1}{0}{1} \suchthat  c_2,c_3\in\Re}
\end{align*}
（读作「$C$~perp」）。
于是我们得到 $\Re^3$ 的如下一组基。
\begin{equation*}
 \sequence{\votepreflist{1}{1}{1},
             \votepreflist{-1}{1}{0},
             \votepreflist{-1}{0}{1}}
\tag*{}\end{equation*}
我们可以关于这组基来表示选票，
从而把选票分解成一个循环部分和一个非循环部分。
\textit{（给学过本章选读小节的读者的提示：~也就是说，
该空间是 $C$ 与~$C^\perp$ 的直和。）}

例如，考虑上面讨论过的 $D>R>T$ 选民。
我们关于这组基来表示这位选民，
\begin{equation*}
  \begin{linsys}{3}
    c_1  &-  &c_2  &-  &c_3  &=  &-1 \\
    c_1  &+  &c_2  &   &     &=  &1  \\
    c_1  &   &     &+  &c_3  &=  &1  
  \end{linsys}
  \grstep[-\rho_1+\rho_3]{-\rho_1+\rho_2}
  \repeatedgrstep{(-1/2)\rho_2+\rho_3}
  \begin{linsys}{3}
    c_1  &-  &c_2  &-  &c_3      &=  &-1 \\
         &   &2c_2 &+  &c_3      &=  &2  \\
         &   &     &   &(3/2)c_3 &=  &1  
  \end{linsys}
\end{equation*}
得到坐标 $c_1=1/3$、$c_2=2/3$ 与 $c_3=2/3$。
于是
\begin{equation*}
  \votepreflist{-1}{1}{1}
  =\frac{1}{3}\cdot\votepreflist{1}{1}{1}
    +\frac{2}{3}\cdot\votepreflist{-1}{1}{0}
    +\frac{2}{3}\cdot\votepreflist{-1}{0}{1}
  =\votepreflist{1/3}{1/3}{1/3}
   +\votepreflist{-4/3}{2/3}{2/3}
\tag*{}\end{equation*}
就给出了所需的分解，
即分成一个循环部分和一个非循环部分。
\begin{equation*}
  \votinggraphic{3}
=
\votinggraphic{4}
\>+\>
\votinggraphic{5}
\end{equation*}
这样我们看到，这位 $D>R>T$ 选民的
理性偏好列表确实含有一个循环部分。

$T>R>D$ 的选民与刚才讨论的选民正好相反，
因为 `$>$' 符号的方向颠倒了。
这位选民的分解
\begin{equation*}
\votinggraphic{6}
=
\votinggraphic{7}
\>+\>
\votinggraphic{8}
\end{equation*}
表明这对相反的偏好具有恰好相反的分解。
我们说第一位选民具有正的
\definend{旋向}\index{spin}\index{voting paradox!spin}，
因为其循环部分沿着我们为箭头选定的方向，
而第二位选民的旋向是负的。

这两类相反的选民相互抵消这一事实，
反映在他们的选票向量之和为零上。
这提示了另一种统计选举结果的办法。
我们可以先尽可能多地消去相反的偏好列表，
然后把余下的列表相加来确定结果。

下面的表格包含三对相反的偏好列表。
例如，最上面一行包含上面讨论过的选民。
\begin{center}
  \def\betweenrowvspace{5.5ex}
  \makebox[\hsize][l]{%  box runs into margin; suppress overfull warn
  \begin{tabular}{@{}c@{\hspace{0em}}c@{}} 
    % \hline
    \multicolumn{1}{c}{\textit{正旋向}}
    &\multicolumn{1}{c}{\textit{负旋向}}  
     \\ \hline
    {\renewcommand{\arraystretch}{1.4}%%
     \begin{tabular}{@{}c}
      $\text{民主党}>\text{共和党}>\text{第三党}$                     \\
      $\votinggraphic{9} % \votecycle{1}{1}{-1}{3}
      =\votinggraphic{10} %\votecycle{1/3}{1/3}{1/3}{3}
      \>+\>
      \votinggraphic{11}$  %\votecycle{2/3}{2/3}{-4/3}{3}\;{}
    \end{tabular}}
    &{\renewcommand{\arraystretch}{1.4}%%
     \begin{tabular}{c@{}}
      $\text{第三党}>\text{共和党}>\text{民主党}$                     \\
      $\votinggraphic{12}  %\votecycle{-1}{-1}{1}{3}
      =\votinggraphic{13}  %\votecycle{-1/3}{-1/3}{-1/3}{3}
      \>+\>
      \votinggraphic{14}$  %\votecycle{-2/3}{-2/3}{4/3}{3}\;{}
    \end{tabular}}                  \\[\betweenrowvspace] % \hline
    {\renewcommand{\arraystretch}{1.4}%%
    \begin{tabular}{@{}c}
      $\text{共和党}>\text{第三党}>\text{民主党}$                     \\
      $\votinggraphic{15}  %\votecycle{-1}{1}{1}{3}
      =\votinggraphic{16} %\votecycle{1/3}{1/3}{1/3}{3}
      \>+\>
      \votinggraphic{17}$  %\votecycle{-4/3}{2/3}{2/3}{3}\;{}
    \end{tabular}}
    &{\renewcommand{\arraystretch}{1.4}%%
     \begin{tabular}{c@{}}
      $\text{民主党}>\text{第三党}>\text{共和党}$                     \\
      $\votinggraphic{18}  %\votecycle{1}{-1}{-1}{3}
      =\votinggraphic{19}  %\votecycle{-1/3}{-1/3}{-1/3}{3}
      \>+\>
      \votinggraphic{20}$  %\votecycle{4/3}{-2/3}{-2/3}{3}\;{}
    \end{tabular}}                  \\[\betweenrowvspace] % \hline
    {\renewcommand{\arraystretch}{1.4}%%
    \begin{tabular}{@{}c}
      $\text{第三党}>\text{民主党}>\text{共和党}$                     \\
      $\votinggraphic{21}  %\votecycle{1}{-1}{1}{3}
      =\votinggraphic{22}  %\votecycle{1/3}{1/3}{1/3}{3}
      \>+\>
      \votinggraphic{23}$  %\votecycle{2/3}{-4/3}{2/3}{3}\;{}
    \end{tabular}}
    &{\renewcommand{\arraystretch}{1.4}%%
     \begin{tabular}{c@{}}
      $\text{共和党}>\text{民主党}>\text{第三党}$                     \\
      $\votinggraphic{24}  %\votecycle{-1}{1}{-1}{3}
      =\votinggraphic{25}  %\votecycle{-1/3}{-1/3}{-1/3}{3}
      \>+\>
      \votinggraphic{26}$  %\votecycle{-2/3}{4/3}{-2/3}{3}\;{}
    \end{tabular}}                  % \\ \hline
  \end{tabular}
  }
  \label{table:Voting}
\end{center}
如果我们按刚才描述的方式进行选举，那么在尽可能多地
消去相反的选民对之后，
将剩下三组偏好列表：~一组来自第一行，
一组来自第二行，一组来自第三行。
我们将以证明下面这一点来结束：
只有当这三组的旋向都相同时，
投票悖论才可能发生。
也就是说，要出现投票悖论，
剩下的三组必须全部来自表格的左侧，
或全部来自右侧
（见 \nearbyexercise{exer:CancelPolSci}）。
这表明多数循环与我们正在使用的这种分解之间存在某种联系\Dash 
投票悖论只有在循环偏好的倾向相互加强时才可能发生。

为给出证明，假设我们已经消去了相反的偏好顺序，
三行中的每一行都剩下一组偏好列表。
考虑第一行剩下的偏好列表。
它们可能来自第一行的左侧或右侧 
（或者介于两者之间，因为这些列表可能恰好完全抵消）。
我们记
\begin{equation*}
  \votinggraphic{27}
\end{equation*}
其中 $a$ 是一个整数：如果剩下的列表在左侧，它为正；
如果列表在右侧，它为负；
如果恰好完全抵消，则为零。
类似地，第二行和第三行各有一个整数 $b$ 和 $c$，
它们各自可正、可负、可为零。

那么这次选举就由下述和式决定。
\begin{equation*}
  \votinggraphic{27}  %\votecycle{a}{a}{-a}{4} 
  +  
  \votinggraphic{28}  %\votecycle{-b}{b}{b}{4} 
  +  
  \votinggraphic{29}  %\votecycle{c}{-c}{c}{4}  
  =\hspace{.30in}
  \votinggraphic{30}  %\votecycle{a-b+c}{a+b-c}{-a+b+c}{4}
  \hspace{.30in}\hbox{}
\end{equation*}
当右边总循环中的三个数，
$-a+b+c$、$a-b+c$ 与 $a+b-c$，
全部非负或全部非正时，就会出现投票悖论。
我们将证明，这种情况只在
$a$、$b$ 与~$c$ 三个数全部非负，
或三个数全部非正时才会发生。

设总循环数非负；另一种情形类似。
\begin{equation*}
  \begin{linsys}{3}
    -a &+ &b &+ &c &\geq &0 \\
     a &- &b &+ &c &\geq &0 \\
     a &+ &b &- &c &\geq &0 
  \end{linsys}
\end{equation*}
把前两行相加，得 $c\geq 0$。
把第一行与第三行相加，得 $b\geq 0$。
而第二行与第三行相加，得 $a\geq 0$。
因此，如果总循环非负，那么每一行剩下的
偏好列表都来自表格的左侧。
证明完毕。

这个结果只是说，三组旋向相同是出现多数循环的
必要条件。
它并不充分；见
\nearbyexercise{exer:VoterCondNotSuff}。

投票理论及相关课题是目前研究的对象。
有许多引人入胜的结果，特别是
K~Arrow \cite{Arrow} 得到的那些，
他因这项工作的一部分而获得诺贝尔奖，
其结果表明没有任何投票制度是完全公平的
（就「公平」的一个合理解义而言）。
一些好的入门文章有 \cite{Gardner70}、
\cite{Gardner74}、\cite{Gardner80} 与 \cite{NeimiRiker}。
\cite{Taylor} 是一本易读的教材。
\cite{GamingVote} 中列出的来自美国近期政治史的
一长串案例表明，这些悖论在实践中经常被操纵。
（另一方面，相当近期的研究表明，
计算如何操纵一场选举在一般情况下可能是不可行的，
但这超出了我们的范围。）
本专题取材于 \cite{Zwicker}。
\emph{（作者注：~我要感谢~Zwicker 教授
热心而富有启发性的讨论。）}

\begin{exercises}
  \item \label{IrrIndVote}
    这里给出一种选民可能具有循环偏好的合理方式。
    假设这位选民按照三条标准中的每一条
    给每位候选人排序。
    \begin{exparts}
      \partsitem 列出一张表，行标记为「民主党」、
         「共和党」和「第三党」，列标记为
         「品格」、「经验」和「政策」。
         在每一列中，把某位候选人
         排在最偏好，
         把另一位排在中间，把剩下的一位
         排在最不偏好。
      \partsitem 在这个排序中，民主党是否在三条标准中的
         （至少）两条上偏好于共和党，
         还是相反？
         共和党是否偏好于第三党？
      \partsitem 刚刚构造的表是否具有
         循环的偏好顺序？
         如果没有，请构造一个具有循环偏好顺序的表。
    \end{exparts}
    因此，选民有可能在诸位候选人之间具有循环偏好。
    然而，上面描述的悖论是：即使每位选民都具有
    直线式的偏好列表，
    整个群体仍可能出现循环偏好。
    \begin{answer}
      这个例子给出单个选民的一个非理性偏好顺序。
      \smallskip
      \begin{center}
        \begin{tabular}{l|c|c|c|}
             \multicolumn{1}{c}{\ }
                  &\multicolumn{1}{c}{\textit{品格}}
                  &\multicolumn{1}{c}{\textit{经验}}
                  &\multicolumn{1}{c}{\textit{政策}}   \\ 
          \cline{2-4}
         \textit{民主党}   &最偏好    &中间  &最不偏好   \\
         \textit{共和党} &中间  &最不偏好   &最偏好    \\
         \textit{第三党}      &最不偏好   &最偏好    &中间  \\
          \cline{2-4}
        \end{tabular}
      \end{center}
      \smallskip
      在品格和经验上，民主党优于共和党。
      在品格和政策上，共和党胜过第三党。
      而在经验和政策上，第三党击败民主党。
    \end{answer}
  \item \label{VoteCalcTable} 
    计算分解表中的各个值。
    \begin{answer}
      首先，把本专题中讲过的 $D>R>T$ 分解
      \begin{equation*}
        \votepreflist{-1}{1}{1}
          =\frac{1}{3}\cdot\votepreflist{1}{1}{1}
           +\frac{2}{3}\cdot\votepreflist{-1}{1}{0}
           +\frac{2}{3}\cdot\votepreflist{-1}{0}{1}
      \end{equation*}
      与相反的 $T>R>D$ 选民的分解
      \begin{equation*}
        \votepreflist{1}{-1}{-1}
          =d_1\cdot\votepreflist{1}{1}{1}
           +d_2\cdot\votepreflist{-1}{1}{0}
           +d_3\cdot\votepreflist{-1}{0}{1}
      \tag*{}\end{equation*}
      相比较。
      显然，第二个是第一个的相反向量，因此
      $d_1=-1/3$、$d_2=-2/3$、$d_3=-2/3$。
      这一原则对任何一对相反的选民都成立，所以我们只需
      对来自第二行的一位选民与来自第三行的
      一位选民作计算。
      对于第二行中具有正旋向的选民，
      \begin{equation*}
        \begin{linsys}{3}
          c_1  &-  &c_2  &-  &c_3  &=  &1 \\
          c_1  &+  &c_2  &   &     &=  &1  \\
          c_1  &   &     &+  &c_3  &=  &-1  
        \end{linsys}
        \grstep[-\rho_1+\rho_3]{-\rho_1+\rho_2}
        \repeatedgrstep{(-1/2)\rho_2+\rho_3}
        \begin{linsys}{3}
          c_1  &-  &c_2  &-  &c_3      &=  &1 \\
               &   &2c_2 &+  &c_3      &=  &0  \\
               &   &     &   &(3/2)c_3 &=  &-2  
        \end{linsys}
      \tag*{}\end{equation*}
      得 $c_3=-4/3$、$c_2=2/3$、$c_1=1/3$。
      对于第三行中具有正旋向的选民，
      \begin{equation*}
        \begin{linsys}{3}
          c_1  &-  &c_2  &-  &c_3  &=  &1 \\
          c_1  &+  &c_2  &   &     &=  &-1  \\
          c_1  &   &     &+  &c_3  &=  &1  
        \end{linsys}
        \grstep[-\rho_1+\rho_3]{-\rho_1+\rho_2}
        \repeatedgrstep{(-1/2)\rho_2+\rho_3}
        \begin{linsys}{3}
          c_1  &-  &c_2  &-  &c_3      &=  &1 \\
               &   &2c_2 &+  &c_3      &=  &-2  \\
               &   &     &   &(3/2)c_3 &=  &1  
        \end{linsys}
        \tag*{}\end{equation*}
        得 $c_3=2/3$、$c_2=-4/3$、$c_1=1/3$。
      \end{answer}
  \item \label{exer:CancelPolSci} 
    对政治学课堂的模拟选举执行相反偏好顺序的抵消。
    剩下的偏好是全部来自表格左边的三行，
    还是来自右边？
    \begin{answer}
      模拟选举按第一个表所示的方式对应于
      第~\pageref{table:Voting}~页的表格，
      抵消之后的结果是第二个表。  
     \smallskip
      \begin{center}
         % \renewcommand{\arraystretch}{1.2}
         \begin{tabular}{|c|c|} \hline
           \textit{正旋向}
           &\textit{负旋向}  \\ \hline
           \begin{tabular}{l}
             $D>R>T$        \\
             $5$ 名选民
           \end{tabular}
          &\begin{tabular}{l}
             $T>R>D$          \\
             $2$ 名选民
           \end{tabular}                  \\ \hline
          \begin{tabular}{l}
             $R>T>D$           \\
             $8$ 名选民
          \end{tabular}
         &\begin{tabular}{l}
             $D>T>R$             \\
             $4$ 名选民
           \end{tabular}                  \\ \hline
          \begin{tabular}{l}
            $T>D>R$             \\
            $8$ 名选民
          \end{tabular}
         &\begin{tabular}{l}
            $R>D>T$            \\
            $2$ 名选民
          \end{tabular}                  \\ \hline
        \end{tabular}
        \qquad
         \begin{tabular}{|c|c|} \hline
           \textit{正旋向}
           &\textit{负旋向}  \\ \hline
           \begin{tabular}{l}
             $D>R>T$        \\
             $3$ 名选民
           \end{tabular}
          &\begin{tabular}{l}
             $T>R>D$          \\
             \textit{-已抵消-}
           \end{tabular}                  \\ \hline
          \begin{tabular}{l}
             $R>T>D$           \\
             $4$ 名选民
          \end{tabular}
         &\begin{tabular}{l}
             $D>T>R$             \\
             \textit{-已抵消-}
           \end{tabular}                  \\ \hline
          \begin{tabular}{l}
            $T>D>R$             \\
            $6$ 名选民
          \end{tabular}
         &\begin{tabular}{l}
            $R>D>T$            \\
            \textit{-已抵消-}
          \end{tabular}                  \\ \hline
        \end{tabular}
       \end{center}
       \smallskip
       三组全部来自表格的同一侧（左侧），
       正如本专题的结果所指出的必然情形。
       现在我们可以统计，利用抵消后的数字
      \begin{equation*}
        % \psset{xunit=4pt,yunit=4pt,runit=4pt} 
        3\cdot \votinggraphic{31}
%          \pspicture[.4](-10.5,-4.2)(10,4) %\psgrid(-10,-10)(10,10)
%            \SpecialCoor
%            \rput(3;90){\small $D$}
%            \psarc{->}(0,0){3}{10}{70} %-rd->
%               \rput[lb](3.3;20){\scriptsize $1$ voter} 
%            \rput[B](3;210){\small $T$}
%            \psarc{->}(0,0){3}{220}{320} %-tr->
%               \rput[t](3.1;270){\scriptsize $1$ voter} 
%            \rput[B](3;330){\small $R$}
%            \psarc{->}(0,0){3}{110}{175} %-dt->
%               \rput[rb](3.3;160){\scriptsize $-1$ voter} 
%          \endpspicture
        +4\cdot\votinggraphic{32}
%          \pspicture[.4](-10,-4.2)(10,4) %\psgrid(-10,-10)(10,10)
%            \SpecialCoor
%            \rput(3;90){\small $D$}
%            \psarc{->}(0,0){3}{10}{70} %-rd->
%               \rput[lb](3.3;20){\scriptsize $-1$ voter} 
%            \rput[B](3;210){\small $T$}
%            \psarc{->}(0,0){3}{220}{320} %-tr->
%               \rput[t](3.1;270){\scriptsize $1$ voter} 
%            \rput[B](3;330){\small $R$}
%            \psarc{->}(0,0){3}{110}{175} %-dt->
%               \rput[rb](3.3;160){\scriptsize $1$ voter} 
%          \endpspicture
        +6\cdot\votinggraphic{33}
%          \pspicture[.4](-9,-4.2)(10,4) %\psgrid(-10,-10)(10,10)
%            \SpecialCoor
%            \rput(3;90){\small $D$}
%            \psarc{->}(0,0){3}{10}{70} %-rd->
%               \rput[lb](3.3;20){\scriptsize $1$ voter} 
%            \rput[B](3;210){\small $T$}
%            \psarc{->}(0,0){3}{220}{320} %-tr->
%               \rput[t](3.1;270){\scriptsize $-1$ voter} 
%            \rput[B](3;330){\small $R$}
%            \psarc{->}(0,0){3}{110}{175} %-dt->
%               \rput[rb](3.3;160){\scriptsize $1$ voter} 
%          \endpspicture
        = \votinggraphic{34}
%          \pspicture[.4](-9,-4.2)(10,4) %\psgrid(-10,-10)(10,10)
%            \SpecialCoor
%            \rput(3;90){\small $D$}
%            \psarc{->}(0,0){3}{10}{70} %-rd->
%               \rput[lb](3.3;20){\scriptsize $5$ voters} 
%            \rput[B](3;210){\small $T$}
%            \psarc{->}(0,0){3}{220}{320} %-tr->
%               \rput[t](3.1;270){\scriptsize $1$ voter} 
%            \rput[B](3;330){\small $R$}
%            \psarc{->}(0,0){3}{110}{175} %-dt->
%               \rput[rb](3.3;160){\scriptsize $7$ voters} 
%          \endpspicture
      \tag*{}\end{equation*}
       得到同样的结果。
    \end{answer}
  \item \label{exer:VoterCondNotSuff} 
    必要条件——只有当抵消后剩下的三个偏好列表
    具有相同旋向时才可能出现投票悖论——
    并不充分。 
    \begin{exparts}
      % \partsitem Continuing the nonnegative total cycle case 
      %   considered in the proof, 
      %   use the two inequalities $0\leq a-b+c$ and $0\leq -a+b+c$
      %   to show that $\absval{a-b}\leq c$.
      % \partsitem Also show that $c\leq a+b$, and hence that
      %   $\absval{a-b}\leq c\leq a+b$. 
      \partsitem 给出一个选举的例子：其中存在多数循环，
        而添加一名具有相同旋向的选民
        会使循环消失。
      \partsitem 反过来是否可能？添加一名具有
        「错误」旋向的选民是否会使循环出现？
      \partsitem 给出一个得到多数循环的
        充要条件。
    \end{exparts}
    \begin{answer}
      \begin{exparts}
        % \partsitem We can rewrite the two as $-c\leq a-b$ and $-c\leq b-a$.
        %   Either $a-b$ or $b-a$ is nonpositive and so $-c\leq-\absval{a-b}$,
        %   as required.
        % \partsitem This is immediate from the supposition that $0\leq a+b-c$.
        \partsitem 一个平凡的例子从零选民的选举开始，
          它有一个平凡的多数循环，然后添加任意一名选民。
          一个更有趣的例子是取政治学课堂的模拟
          选举并添加两名 $T>D>R$ 选民
          （为了满足题目中
          「再添加一名选民」的要求，
          我们可以一次添加一名）。
          新选民的旋向为正，这正是原模拟选举中
          抵消后剩下的选票的旋向。
          下面是所得的选民表，旁边是
          抵消的结果。
      \smallskip
      \begin{center}
         % \renewcommand{\arraystretch}{1.3}
         \begin{tabular}{|c|c|} \hline
           \textit{正旋向}
           &\textit{负旋向}  \\ \hline
           \begin{tabular}{l}
             $D>R>T$        \\
             $5$ 名选民
           \end{tabular}
          &\begin{tabular}{l}
             $T>R>D$          \\
             $2$ 名选民
           \end{tabular}                  \\ \hline
          \begin{tabular}{l}
             $R>T>D$           \\
             $8$ 名选民
          \end{tabular}
         &\begin{tabular}{l}
             $D>T>R$             \\
             $4$ 名选民
           \end{tabular}                  \\ \hline
          \begin{tabular}{l}
            $T>D>R$             \\
            $10$ 名选民
          \end{tabular}
         &\begin{tabular}{l}
            $R>D>T$            \\
            $2$ 名选民
          \end{tabular}                  \\ \hline
        \end{tabular}
        \qquad
         \begin{tabular}{|c|c|} \hline
           \textit{正旋向}
           &\textit{负旋向}  \\ \hline
           \begin{tabular}{l}
             $D>R>T$        \\
             $3$ 名选民
           \end{tabular}
          &\begin{tabular}{l}
             $T>R>D$          \\
             \textit{-已抵消}
           \end{tabular}                  \\ \hline
          \begin{tabular}{l}
             $R>T>D$           \\
             $4$ 名选民
          \end{tabular}
         &\begin{tabular}{l}
             $D>T>R$             \\
             \textit{-已抵消-}
           \end{tabular}                  \\ \hline
          \begin{tabular}{l}
            $T>D>R$             \\
            $8$ 名选民
          \end{tabular}
         &\begin{tabular}{l}
            $R>D>T$            \\
            \textit{-已抵消-}
          \end{tabular}                  \\ \hline
        \end{tabular}
       \end{center}
       \smallskip
       这是利用抵消后的数字所作的选举。
       \begin{equation*}
         % \psset{xunit=4pt,yunit=4pt,runit=4pt} 
         3\cdot\votinggraphic{35}
%          \pspicture[.4](-10.5,-4.2)(10,4) %\psgrid(-10,-10)(10,10)
%            \SpecialCoor
%            \rput(3;90){\small $D$}
%            \psarc{->}(0,0){3}{10}{70} %-rd->
%               \rput[lb](3.3;20){\scriptsize $1$ voter} 
%            \rput[B](3;210){\small $T$}
%            \psarc{->}(0,0){3}{220}{320} %-tr->
%               \rput[t](3.1;270){\scriptsize $1$ voter} 
%            \rput[B](3;330){\small $R$}
%            \psarc{->}(0,0){3}{110}{175} %-dt->
%               \rput[rb](3.3;160){\scriptsize $-1$ voter} 
%          \endpspicture
         +4\cdot\votinggraphic{36}
%          \pspicture[.4](-10,-4.2)(10,4) %\psgrid(-10,-10)(10,10)
%            \SpecialCoor
%            \rput(3;90){\small $D$}
%            \psarc{->}(0,0){3}{10}{70} %-rd->
%               \rput[lb](3.3;20){\scriptsize $-1$ voter} 
%            \rput[B](3;210){\small $T$}
%            \psarc{->}(0,0){3}{220}{320} %-tr->
%               \rput[t](3.1;270){\scriptsize $1$ voter} 
%            \rput[B](3;330){\small $R$}
%            \psarc{->}(0,0){3}{110}{175} %-dt->
%               \rput[rb](3.3;160){\scriptsize $1$ voter} 
%          \endpspicture
         +8\cdot\votinggraphic{37}
%          \pspicture[.4](-9,-4.2)(10,4) %\psgrid(-10,-10)(10,10)
%            \SpecialCoor
%            \rput(3;90){\small $D$}
%            \psarc{->}(0,0){3}{10}{70} %-rd->
%               \rput[lb](3.3;20){\scriptsize $1$ voter} 
%            \rput[B](3;210){\small $T$}
%            \psarc{->}(0,0){3}{220}{320} %-tr->
%               \rput[t](3.1;270){\scriptsize $-1$ voter} 
%            \rput[B](3;330){\small $R$}
%            \psarc{->}(0,0){3}{110}{175} %-dt->
%               \rput[rb](3.3;160){\scriptsize $1$ voter} 
%          \endpspicture
         =\votinggraphic{38}
%          \pspicture[.4](-9,-4.2)(10,4) %\psgrid(-10,-10)(10,10)
%            \SpecialCoor
%            \rput(3;90){\small $D$}
%            \psarc{->}(0,0){3}{10}{70} %-rd->
%               \rput[lb](3.3;20){\scriptsize $7$ voters} 
%            \rput[B](3;210){\small $T$}
%            \psarc{->}(0,0){3}{220}{320} %-tr->
%               \rput[t](3.1;270){\scriptsize $-1$ voter} 
%            \rput[B](3;330){\small $R$}
%            \psarc{->}(0,0){3}{110}{175} %-dt->
%               \rput[rb](3.3;160){\scriptsize $9$ voters} 
%          \endpspicture
      \tag*{}\end{equation*}
       多数循环确实消失了。
      \partsitem 把前面的练习反过来。
         也就是说，在前一个练习的结果上添加一名选民，
         把使循环消失的那位选民抵消掉。 
      \partsitem 一个这样的条件是：抵消之后，
         三个数全部非负或全部非正，
         并且：~$\absval{c}<\absval{a+b}$、$\absval{b}<\absval{a+c}$ 
         与 $\absval{a}<\absval{b+c}$。
         这由下图得出。 
         \begin{equation*}
           \votinggraphic{27}  %\votecycle{a}{a}{-a}{4} 
           +  
           \votinggraphic{28}  %\votecycle{-b}{b}{b}{4} 
           +  
           \votinggraphic{29}  %\votecycle{c}{-c}{c}{4}  
           =\hspace{.30in}
           \votinggraphic{30}  %\votecycle{a-b+c}{a+b-c}{-a+b+c}{4}
           \hspace{.30in}\hbox{}  
         \end{equation*}
%          \begin{center}
%            \votecycle{a}{a}{-a}{4} 
%            +  
%            \votecycle{-b}{b}{b}{4} 
%            +  
%            \votecycle{c}{-c}{c}{4}  
%            =\hspace{.30in}  
%            \votecycle{a-b+c}{a+b-c}{-a+b+c}{4}\hspace{.30in}\hbox{}  
%          \end{center}
      \end{exparts}
    \end{answer}
\item 一名选民的选举不可能出现多数循环，
    因为我们已施加了选民列表必须理性的要求。
    \begin{exparts}
      \partsitem 证明两名选民的选举可能出现多数循环。
        （如果三个群体总数全部非负或全部非正——也就是说，
        我们允许群体偏好中出现一些零——
        我们就把群体偏好视为多数循环。）
      \partsitem 证明对于任何大于一的选民数，
        都存在涉及该数目选民的一场选举，
        其结果是一个多数循环。
    \end{exparts}
    \begin{answer}
      \begin{exparts}
        \partsitem 两名选民的选举可能以两种方式出现多数循环。
          第一，两名选民可能恰好相反，
          抵消后得到平凡的选举（多数循环全为零）。
          第二，两名选民可能具有相同旋向但来自
          不同的行，如下所示。
          \begin{equation*}
            % \psset{xunit=4pt,yunit=4pt,runit=4pt} 
            1\cdot \votinggraphic{39}
%             \pspicture[.4](-10.5,-4.2)(10,4) %\psgrid(-10,-10)(10,10)
%               \SpecialCoor
%               \rput(3;90){\small $D$}
%               \psarc{->}(0,0){3}{10}{70} %-rd->
%                  \rput[lb](3.3;20){\scriptsize $1$ voter} 
%               \rput[B](3;210){\small $T$}
%               \psarc{->}(0,0){3}{220}{320} %-tr->
%                  \rput[t](3.1;270){\scriptsize $1$ voter} 
%               \rput[B](3;330){\small $R$}
%               \psarc{->}(0,0){3}{110}{175} %-dt->
%                  \rput[rb](3.3;160){\scriptsize $-1$ voter} 
%             \endpspicture
            +1\cdot\votinggraphic{40}
%             \pspicture[.4](-10,-4.2)(10,4) %\psgrid(-10,-10)(10,10)
%               \SpecialCoor
%               \rput(3;90){\small $D$}
%               \psarc{->}(0,0){3}{10}{70} %-rd->
%                  \rput[lb](3.3;20){\scriptsize $-1$ voter} 
%               \rput[B](3;210){\small $T$}
%               \psarc{->}(0,0){3}{220}{320} %-tr->
%                  \rput[t](3.1;270){\scriptsize $1$ voter} 
%               \rput[B](3;330){\small $R$}
%               \psarc{->}(0,0){3}{110}{175} %-dt->
%                  \rput[rb](3.3;160){\scriptsize $1$ voter} 
%             \endpspicture
            +0\cdot\votinggraphic{41}
%             \pspicture[.4](-9,-4.2)(10,4) %\psgrid(-10,-10)(10,10)
%               \SpecialCoor
%               \rput(3;90){\small $D$}
%               \psarc{->}(0,0){3}{10}{70} %-rd->
%                  \rput[lb](3.3;20){\scriptsize $1$ voter} 
%               \rput[B](3;210){\small $T$}
%               \psarc{->}(0,0){3}{220}{320} %-tr->
%                  \rput[t](3.1;270){\scriptsize $-1$ voter} 
%               \rput[B](3;330){\small $R$}
%               \psarc{->}(0,0){3}{110}{175} %-dt->
%                  \rput[rb](3.3;160){\scriptsize $1$ voter} 
%             \endpspicture
            =\votinggraphic{42}
%             \pspicture[.4](-9,-4.2)(10,4) %\psgrid(-10,-10)(10,10)
%               \SpecialCoor
%               \rput(3;90){\small $D$}
%               \psarc{->}(0,0){3}{10}{70} %-rd->
%                  \rput[lb](3.3;20){\scriptsize $0$ voters} 
%               \rput[B](3;210){\small $T$}
%               \psarc{->}(0,0){3}{220}{320} %-tr->
%                  \rput[t](3.1;270){\scriptsize $2$ voters} 
%               \rput[B](3;330){\small $R$}
%               \psarc{->}(0,0){3}{110}{175} %-dt->
%                  \rput[rb](3.3;160){\scriptsize $0$ voters} 
%             \endpspicture
          \tag*{}\end{equation*}
        \partsitem 有两种情形。
          偶数个选民可以平分成相反的两半，
          例如一半选民是 $D>R>T$，一半是 $T>R>D$。
          那么抵消给出平凡的选举。
          如果选民数大于一且为奇数（形如
          $2k+1$，其中 $k>0$），那么利用证明中的循环图，
          \begin{equation*}
            \votinggraphic{27}  %\votecycle{a}{a}{-a}{4} 
            +  
            \votinggraphic{28}  %\votecycle{-b}{b}{b}{4} 
            +  
            \votinggraphic{29}  %\votecycle{c}{-c}{c}{4}  
            =\hspace{.30in}
            \votinggraphic{30}  %\votecycle{a-b+c}{a+b-c}{-a+b+c}{4}
            \hspace{.30in}\hbox{}  
          \end{equation*}
%          \begin{center}
%            \votecycle{a}{a}{-a}{4} 
%            +  
%            \votecycle{-b}{b}{b}{4} 
%            +  
%            \votecycle{c}{-c}{c}{4}  
%            =\hspace{.30in}  
%            \votecycle{a-b+c}{a+b-c}{-a+b+c}{4}\hspace{.30in}\hbox{}  
%          \end{center}
          我们可以取 $a=k$、$b=k$、$c=1$。
          因为 $k>0$，这就是一个多数循环。
       \end{exparts}
    \end{answer}
\item \label{exer:PerpIsSubSp} 
    设 $U$ 是 $\Re^3$ 的一个子空间。
    证明与 $U$ 中每个向量都垂直的向量所成的集合
    $U^\perp=\set{\vec{v}\suchthat 
              \text{$\vec{v}\dotprod\vec{u}=0$ 对所有 $\vec{u}\in U$}}$
    也是 $\Re^3$ 的子空间。
    若 $U$ 不是子空间，这是否仍然成立？
    \begin{answer}
      它非空，因为它包含零向量。
      为证它对其两个成员的线性组合封闭，
      设 $\vec{v}_1$ 与 $\vec{v}_2$ 属于 $U^\perp$，
      考虑 $c_1\vec{v}_1+c_2\vec{v}_2$。
      对任意 $\vec{u}\in U$，
      \begin{equation*}
        (c_1\vec{v}_1+c_2\vec{v}_2)\dotprod\vec{u}
        =c_1(\vec{v}_1\dotprod\vec{u})+c_2(\vec{v}_2\dotprod\vec{u})
        =c_1\cdot 0+c_2\cdot 0=0         
      \tag*{}\end{equation*}
      所以 $c_1\vec{v}_1+c_2\vec{v}_2\in U^\perp$。

      至于 $U$ 是一个子集但不是子空间时是否仍然成立，
      答案是肯定的。
    \end{answer}
\end{exercises}
\index{voting paradox|)}
\endinput


